\section{Execution Model}
\label{sec:execution}

\subsection{Two levels of resumption}

A 275,218-call study is not a single job; it is 38 stages, each of which may be interrupted. The harness resumes at two granularities.

\textbf{Stage level.} A state file records completed stage keys of the form \texttt{layer/model/stage}, plus a separate list of contract-exempt stages. A completed stage is skipped on restart; a stage that exits non-zero, or that exits zero with error rows in its output, is \emph{not} marked complete, so a restart retries it. A run-level configuration fingerprint hashes the three manifests and the prompt/schema modules; if the fingerprint differs from the one recorded in the state file, the run refuses to resume and asks for \texttt{-Reset}. The delivered state predates the guard; it now carries a post-hoc recomputation from the same seven files (\texttt{config\_fingerprint\_recomputed}), while the cross-modal and relay arms persist the fingerprint at run time. The delivered run has 38 stage entries: 30 completed, of which 8 are flagged contract-exempt for structured-contract metrics (six Open-i stages and two LLaVA stages). Open-i prose metrics are still computed; only the structured-contract metrics are excluded.

\textbf{Row level.} Each runner writes a JSONL prediction file through \texttt{DurableJsonl}. Every record is appended, flushed, and fsynced; the durable byte offset is written to an atomic sidecar (temp file $\rightarrow$ \texttt{os.replace}). On reopen, the reader truncates the JSONL back to the durable offset, so a torn tail cannot survive a crash, and reports exactly which \texttt{(image\_id, condition, severity)} keys are present. A failed row is retried by default; a deterministic failure can be made terminal with \texttt{retry\_failed=False}.

\textbf{Replacement, not duplication.} Retrying a failed key inside one run must not leave both the failure and the success in the file. \texttt{write\_many} collapses repeats inside the batch, then replaces the existing record in place using an in-memory key$\rightarrow$index map. The original implementation scanned the full record list for every key; we replaced that scan with the index, turning an $O(n^2)$ path into $O(n)$ for a stage with tens of thousands of rows. A failing progress sidecar is surfaced as \texttt{progress\_error} and a stderr warning rather than swallowed: the next resume falls back to the last complete line, which is safe but costs re-work.

\subsection{Crash and resume tests}

The claim ``a power cut costs at most the in-flight batch'' spans process boundaries, the CLI, the writer, and the resume key, so it is tested end-to-end with a stubbed model (no GPU, no weights). The suite runs the real PadChest and Open-i runners, kills them after a chosen batch, appends a torn JSON line, resumes, and checks the file. Table~\ref{tab:resume} summarizes the 40 checks.

\begin{table}[t]
\centering
\small
\caption{Crash/resume suite (CPU-only, stubbed model). All 40 checks pass. The suite is part of CI and runs in seconds.}
\label{tab:resume}
\begin{tabular}{llr}
\toprule
Runner & Check family & Checks \\
\midrule
Open-i (44) & clean pass, schema, summary, sidecar & 8 \\
 & hard failure: non-zero exit, durable prefix, resume & 7 \\
 & crash: partial file, resume, no duplicates, no failed rows & 8 \\
 & torn tail: truncation, no duplicates, completeness & 4 \\
PadChest (43) & clean pass, blank images, summary, schema compatibility & 7 \\
 & crash/resume: non-zero exit, partial file, no duplicates & 6 \\
\midrule
\multicolumn{2}{r}{Total} & \textbf{40} \\
\bottomrule
\end{tabular}
\end{table}

The strongest check is the third resume: after a clean run, a crash, a torn tail, and two resumes, a third invocation must report nothing to do and leave the JSONL \emph{byte-identical}. That is what ``resumable'' has to mean when it is relied on dozens of times across a multi-day run.
